\documentclass[10pt]{article}
\usepackage[utf8]{inputenc}
\usepackage[T1]{fontenc}
\usepackage[a4paper,margin=1.7cm]{geometry}
\usepackage{enumitem}
\usepackage{xcolor}
\usepackage{titlesec}
\usepackage{parskip}

\definecolor{azul}{HTML}{2F5FB3}
\definecolor{gris}{HTML}{555555}

\titleformat{\section}{\large\bfseries\color{azul}}{}{0pt}{}
\titlespacing*{\section}{0pt}{8pt}{2pt}
\setlist{nosep,leftmargin=1.4em}
\pagestyle{empty}

\begin{document}

\begin{center}
    {\LARGE\bfseries\color{azul} Limpiador de bases de contactos}\\[2pt]
    {\small\color{gris} Guía rápida de uso}
\end{center}

\vspace{2pt}

\noindent Esta herramienta revisa una lista de correos y determina cuáles sirven para
enviar campañas, reduciendo los rebotes (correos que no llegan). Acepta archivos
\textbf{Excel} y \textbf{CSV}, y no modifica tu archivo original: genera uno nuevo, ya limpio.

\section{El flujo completo}
\begin{center}
    \textbf{Subir correos} $\rightarrow$ \textbf{Analizar / limpiar} $\rightarrow$
    \textbf{Chequeo SMTP} (opcional) $\rightarrow$ \textbf{Cargar en HubSpot} (opcional)
\end{center}
\noindent Subes tu lista, la herramienta la analiza y clasifica, opcionalmente haces un
chequeo más profundo (SMTP), y opcionalmente marcas el resultado en HubSpot. Solo el
primer paso es obligatorio; los demás, cuando los necesites.

\section{Qué puede y qué no puede hacer}
\textbf{Sí puede:} validar que el correo esté bien escrito, que su dominio pueda recibir
mensajes (registro MX), detectar dominios desechables, cuentas genéricas (\texttt{info@}),
corregir errores de tipeo y marcar duplicados. Con el chequeo SMTP, además intenta
confirmar el buzón contra el servidor de destino.

\textbf{No puede:} saber si una cuenta está \emph{activa} o es ``zombi'', ni detectar
trampas de spam, ni confirmar buzones en dominios que aceptan todo (catch-all). Para ese
nivel de certeza existe software especializado de pago como \textbf{ZeroBounce}. Esta
herramienta cubre bien la mayoría de los rebotes evitables, gratis y con tus propios datos.

\section{Cómo usarla (3 pasos)}
\begin{enumerate}[label=\textbf{\arabic*.}]
    \item \textbf{Sube tu lista.} Arrastra el archivo Excel o CSV. La herramienta detecta
    sola la columna de correos; cámbiala si eligió mal. Aprieta \textbf{Procesar lista}.
    \item \textbf{Revisa el resultado.} Verás cuántos correos son válidos, de riesgo o
    inválidos. Puedes filtrar la tabla por estado o por tipo de error.
    \item \textbf{Descarga la lista limpia.} Elige qué bajar (solo válidos, enviables, o
    todo) y en qué formato (CSV o Excel).
\end{enumerate}

\section{Qué significa cada estado}
\begin{itemize}
    \item \textbf{Válido.} Está bien escrito y su dominio puede recibir correo. Se puede enviar.
    \item \textbf{Riesgo.} Sirve, pero tiene alguna señal a considerar (por ejemplo, es una
    cuenta genérica como \texttt{info@} o \texttt{ventas@}). Enviar con criterio.
    \item \textbf{Inválido.} No conviene enviar: está vacío, mal escrito, el dominio no recibe
    correo o es desechable. Rebotaría.
\end{itemize}

\section{Motivos más frecuentes}
\begin{itemize}
    \item \texttt{sin\_email} --- la fila no traía correo.
    \item \texttt{sintaxis\_invalida} --- el correo está mal escrito.
    \item \texttt{dominio\_sin\_mx} --- el dominio no puede recibir correo (rebota siempre).
    \item \texttt{cuenta\_generica\_rol} --- es una casilla de área (\texttt{info@},
    \texttt{soporte@}), no una persona.
    \item \texttt{dominio\_desechable} --- correo temporal o de un solo uso.
\end{itemize}

\section{Columnas útiles del resultado}
\begin{itemize}
    \item \textbf{estado} y \textbf{motivo} --- el veredicto y su porqué.
    \item \textbf{tipo\_dominio} --- personal (gmail, hotmail) o corporativo.
    \item \textbf{correccion\_typo} --- si se corrigió un error de tipeo (ej. \texttt{gmial.com}
    $\rightarrow$ \texttt{gmail.com}).
    \item \textbf{last\_check\_date} --- cuándo se verificó ese correo.
    \item \textbf{es\_duplicado} --- si el correo está repetido en la lista.
\end{itemize}

\section{Verificación avanzada (SMTP, opcional)}
Es un chequeo más profundo: se conecta al servidor de cada correo (sin enviar nada) para
confirmar mejor si el buzón recibe. Al terminar, \textbf{actualiza el estado de la base}.
Verifica los correos que tengas visibles en la tabla, según tus filtros.
\begin{itemize}
    \item \textbf{Existe} $\rightarrow$ válido. \quad \textbf{No existe / sin MX}
    $\rightarrow$ inválido.
    \item \textbf{Catch-all} (el dominio acepta todo) $\rightarrow$ no cambia el estado: el
    servidor no permite confirmar el buzón.
    \item \textbf{No verificable} (el servidor no responde con certeza) $\rightarrow$ riesgo.
\end{itemize}
\textbf{Ojo:} necesita un correo remitente dedicado (no el corporativo principal) y una red
que no bloquee el puerto 25 --- desde oficinas o casas suele salir ``no verificable'', y eso
es la red, no un error. Está pensado para lotes chicos (máximo 3.000 correos).

\section{Integración con HubSpot (opcional)}
Por defecto no toca HubSpot. Si lo activas, \textbf{marca el estado en cada contacto} (no
borra ni archiva nada). Buscando por correo, escribe tres propiedades:
\begin{itemize}
    \item \textbf{Estado validación email} --- válido / riesgo / inválido.
    \item \textbf{Validado por} --- tu usuario de HubSpot, para que quede quién hizo el chequeo.
    \item \textbf{Última verificación email} --- la fecha.
\end{itemize}
Corre siempre primero en \textbf{Modo simulación} (muestra a cuántos afectaría sin tocar el
CRM). Solo actualiza contactos que ya existen; no crea nuevos.

\section{Consejos}
\begin{itemize}
    \item Para enviar campañas, usa la descarga \textbf{``Solo válidos y sin repetidos''}.
    \item Un correo de riesgo por ser cuenta genérica igual llega; decide según tu campaña.
    \item Las verificaciones se recuerdan: un correo revisado hace poco no se vuelve a chequear.
\end{itemize}

\end{document}
